% =============================================================
%  Chapter 2 — Preliminaries and Related Work
% =============================================================

\chapter{Preliminaries and Related Work}
\label{chap:background}

% ------------------------------------------------------------------
% Macro block — consistent with iclr2026_conference.tex
% ------------------------------------------------------------------
\providecommand{\Ham}{\mathcal{H}}
\providecommand{\Energy}{\mathcal{E}}
\providecommand{\Potential}{\mathcal{R}}
\providecommand{\Constraints}{\mathcal{C}}
\providecommand{\Navigator}{\mathsf{N}}
\providecommand{\PhaseState}{z}
\providecommand{\Pos}{q}
\providecommand{\Mom}{p}
\providecommand{\Mass}{\mathrm{M}}
\providecommand{\J}{\mathbb{J}}
\providecommand{\Lag}{\mathscr{L}}
\providecommand{\mult}{\boldsymbol{\lambda}}
\providecommand{\pen}{\boldsymbol{\rho}}

This chapter has two roles.
The first half (\S\ref{sec:prelim:setting}--\S\ref{sec:prelim:continual})
builds the mathematical vocabulary shared across Chapters~\ref{chap:grlsnam}
and~\ref{chap:material}.
The second half surveys
the literature through the lens established here.

The central thesis lens, stated once and carried throughout:

\begin{quote}\itshape
Rather than learning a fixed policy, reward, or planner, we study
how an agent can learn a \textbf{structured decision dynamics}—a vector field
over phase space whose energy landscape can be \textbf{continually reshaped}
by geometry, material properties, and risk.
\end{quote}

% =======================================================================
\section{Problem Setting and Notation}
\label{sec:prelim:setting}
% =======================================================================

\paragraph{Environment and task.}
The agent operates in a workspace $\mathcal{W} \subset \mathbb{R}^2$ (or $\mathbb{R}^3$)
that is initially unknown.
Obstacles are encoded by an occupancy field
$I : \mathcal{W} \to \{0,1\}$ that is never given in full;
the agent receives only \emph{local} observations through an environment
context $\mathcal{E}_\tau$ at each planning stage $\tau = 0, 1, 2, \ldots$
The task is to navigate from an initial configuration $\vec{x}_0$ to a goal
$\vec{x}_g$ while satisfying geometric and risk constraints, minimizing
both path cost and exposure to latent hazards.

\paragraph{Configuration and phase space.}
The robot's configuration space $\mathcal{Q}$ decomposes as
\begin{equation}
  \mathcal{Q} = \mathcal{Q}_y \times \mathcal{Q}_f \times \mathcal{Q}_o,
  \label{eq:config_decomp}
\end{equation}
where $\vec{y} \in \mathcal{Q}_y$ are sensor variables,
$\vec{c} \in \mathcal{Q}_f \subset \mathbb{R}^2$ is the robot pose in world
coordinates, and $\vec{\psi} \in \mathcal{Q}_o$ collects deformation or
morphology variables.
The full configuration is $\vec{q} = (\vec{y}, \vec{c}, \vec{\psi}) \in
\mathcal{Q}$.
The \textbf{phase space} is the cotangent bundle $T^*\mathcal{Q}$, whose
elements are pairs $(\vec{q}, p)$ with $\vec{q} \in \mathcal{Q}$ and
$p \in T^*_{\vec{q}}\mathcal{Q}$ a covector (generalized momentum dual to
$T_{\vec{q}}\mathcal{Q}$).
We write $\PhaseState = (\vec{q}, \Mom)$ for a generic phase-space point.

\paragraph{Environment context.}
At stage $\tau$, the sensor returns an environment context
$\mathcal{E}_\tau$ containing at minimum:
a local target goal $\vec{c}_{\mathrm{target}}(\mathcal{E}_\tau) \in \mathbb{R}^2$,
an obstacle count $C(\mathcal{E}_\tau) \in \mathbb{N}$,
and signed-clearance functions
$\{d_i(\vec{q};\mathcal{E}_\tau) \geq 0\}_{i=1}^{C(\mathcal{E}_\tau)}$
encoding collision-free space.
In the material-aware stage (Chapter~\ref{chap:material}), $\mathcal{E}_\tau$
additionally contains a material-risk field
$\rho : \mathcal{Q}_f \to [0,1]$ and hazard indicator
$h : \mathcal{Q}_f \to \{0,1\}$.

\paragraph{Notation table.}
Table~\ref{tab:notation} collects the symbols used throughout the thesis.

\begin{table}[h]
\centering
\caption{\textbf{Notation.} Symbols are consistent across all chapters.}
\label{tab:notation}
\small
\setlength{\tabcolsep}{6pt}
\renewcommand{\arraystretch}{1.25}
\begin{tabular}{lll}
\toprule
\textbf{Symbol} & \textbf{Meaning} & \textbf{Space / Type} \\
\midrule
$\mathcal{Q}$               & Configuration space                      & Smooth manifold \\
$T^*\mathcal{Q}$            & Cotangent bundle (phase space)            & Symplectic manifold \\
$\vec{q}$                   & Configuration (pose + sensor + shape)     & $\mathcal{Q}$ \\
$p$                         & Costate / generalized momentum            & $T^*_{\vec{q}}\mathcal{Q}$ \\
$\PhaseState = (\vec{q},p)$ & Phase-space state                         & $T^*\mathcal{Q}$ \\
$H(\vec{q},p;\mathcal{E})$  & Reduced Hamiltonian                       & $C^\infty(T^*\mathcal{Q})$ \\
$\mathcal{R}(\vec{q};\mathcal{E})$ & Potential energy (navigation terms) & $C^\infty(\mathcal{Q})$ \\
$\mathbf{D}(\PhaseState)$   & Damping / dissipation matrix (port-Ham.) & $\mathbb{S}^{2n}_{\succeq 0}$ \\
$\Mass$                     & Mass / mobility matrix                    & $\mathbb{S}^n_{++}$ \\
$\J$                        & Symplectic matrix                         & $\mathbb{R}^{2n\times 2n}$ \\
$\Phi^{\Delta t}_H$         & Time-$\Delta t$ flow map of $H$           & $T^*\mathcal{Q}\to T^*\mathcal{Q}$ \\
$\pi_H(z_t)$                & Policy as phase-space flow (deterministic) & $T^*\mathcal{Q}\to\mathcal{Q}$ \\
$\mathcal{E}_\tau$          & Environment context at stage $\tau$       & Structured record \\
$d_i(\vec{q};\mathcal{E})$  & Signed clearance to obstacle $i$          & $\mathbb{R}$ \\
$C(\mathcal{E})$            & Active obstacle count                     & $\mathbb{N}$ \\
$\alpha_i(\mathcal{E})$     & Barrier weight for obstacle $i$           & $\mathbb{R}_{\geq 0}$ \\
$b_i(s)$                    & Barrier potential applied to $d_i$        & $C^2(\mathbb{R})$ \\
$\rho(\vec{c})$             & Material risk field                       & $[0,1]$ \\
$\xi$                       & Learnable parameters of $\eta_\xi$        & $\mathbb{R}^P$ \\
$\tau$                      & Stage index                               & $\mathbb{N}$ \\
$\mathrm{CVaR}_\beta$       & Conditional Value-at-Risk at level $\beta$ & $\mathbb{R}$ \\
\bottomrule
\end{tabular}
\end{table}

% =======================================================================
\section{Hamiltonian and Port-Hamiltonian Dynamics}
\label{sec:prelim:hamiltonian}
% =======================================================================

\subsection{Classical Hamiltonian Systems}
\label{subsec:prelim:classical_ham}

A \textbf{Hamiltonian system} on $T^*\mathcal{Q} \cong \mathbb{R}^{2n}$
is defined by a smooth scalar $H : T^*\mathcal{Q} \to \mathbb{R}$ called the
\emph{Hamiltonian}, via Hamilton's equations:
\begin{equation}
  \dot{\vec{q}} = \nabla_p H(\vec{q},p), \qquad
  \dot{p}       = -\nabla_{\vec{q}} H(\vec{q},p).
  \label{eq:hamilton_eqs}
\end{equation}
Equivalently, writing $\PhaseState = (\vec{q},p)$ and the standard
symplectic matrix
\begin{equation}
  \J = \begin{pmatrix} 0 & I_n \\ -I_n & 0 \end{pmatrix},
  \label{eq:symplectic_matrix}
\end{equation}
Hamilton's equations take the compact form
$\dot{\PhaseState} = \J\,\nabla_\PhaseState H(\PhaseState)$.

\begin{definition}[Mechanical Hamiltonian]
\label{def:mech_hamiltonian}
A \emph{mechanical Hamiltonian} has the form
\begin{equation}
  H(\vec{q},p;\mathcal{E})
  = \tfrac{1}{2}\,p^\top \Mass^{-1} p
    + \mathcal{R}(\vec{q};\mathcal{E}),
  \label{eq:mech_hamiltonian}
\end{equation}
where $\Mass \in \mathbb{S}^n_{++}$ is a mass/metric matrix and
$\mathcal{R}(\vec{q};\mathcal{E})$ is the \emph{potential energy},
encoding navigation objectives (goal attraction, obstacle repulsion, risk).
The kinetic term $\frac{1}{2}p^\top\Mass^{-1}p$ drives forward momentum;
the gradient $-\nabla_{\vec{q}}\mathcal{R}$ provides the restoring force.
\end{definition}

\paragraph{Flow map and policy.}
The time-$\Delta t$ \emph{flow map} of $H$ advances phase-space state by
integrating \eqref{eq:hamilton_eqs}:
\begin{equation}
  \Phi^{\Delta t}_H : T^*\mathcal{Q} \to T^*\mathcal{Q},
  \qquad
  (\vec{q}_{t+\Delta t}, p_{t+\Delta t})
  = \Phi^{\Delta t}_H(\vec{q}_t, p_t).
  \label{eq:flow_map}
\end{equation}
The \textbf{thesis policy} is the deterministic configuration projection of this flow:
\begin{equation}
  \pi_H(z_t) := \Pi_{\mathcal{Q}}\!\left(\Phi^{\Delta t}_H(z_t)\right),
  \qquad
  z_{t+1} = \Phi^{\Delta t}_H(z_t),
  \label{eq:policy_as_flow}
\end{equation}
where $z_t = (\vec{q}_t, p_t)$ and $\Pi_{\mathcal{Q}}$ projects to configuration.
(The earlier shorthand $\pi(\vec{q}|p)$ used conditional-distribution notation
misleadingly; $\Phi^{\Delta t}_H$ is a \emph{deterministic} map.)
Adapting the policy means reshaping $H$ by adjusting potential coefficients,
adding energy terms and not relearning a flat action mapping.

\paragraph{Symplectic structure.}
A key property of Hamiltonian flows is that they are
\textbf{symplectomorphisms}: they preserve the symplectic 2-form
$\omega = \sum_i dq^i \wedge dp_i$.
This implies volume preservation in phase space (Liouville's theorem) and
controls long-run drift in numerical integration.

\begin{proposition}[Symplectic invariance~\citep{hairer2006geometric}]
\label{prop:symplectic}
Let $H \in C^2(T^*\mathcal{Q})$ be a \emph{conservative} Hamiltonian
(no dissipation or port inputs). Then for all $t$, the map
$\Phi^t_H$ is a symplectomorphism: $(\Phi^t_H)^* \omega = \omega$.
Symplectic integrators~\citep{reich1996symplectic} preserve the discrete
symplectic form exactly (up to roundoff) for separable conservative
Hamiltonians, while generic Euler steps accumulate $O(t)$ drift.
\emph{Note:} once dissipation ($\mathbf{D} \succ 0$) or port inputs ($u$)
are added via the port-Hamiltonian extension~\eqref{eq:port_ham}, the
flow is no longer symplectic.
The dissipative and port-input substeps are handled separately through the
port-Hamiltonian energy balance~\eqref{eq:energy_balance} rather than via
symplectic preservation.
\end{proposition}

\subsection{Port-Hamiltonian Extension}
\label{subsec:prelim:portham}

Classical Hamiltonian systems are conservative; real navigation involves
dissipation (actuator drag, damping) and external forcing (control inputs,
sensor-injected corrections).
The \textbf{port-Hamiltonian framework}~\citep{desai2021porthamiltonian}
extends \eqref{eq:hamilton_eqs} to:
\begin{equation}
  \dot{\PhaseState}
  = \bigl(\J - \mathbf{D}(\PhaseState)\bigr)\,\nabla_\PhaseState H(\PhaseState)
    + \mathbf{B}(\PhaseState)\,u,
  \label{eq:port_ham}
\end{equation}
where:
\begin{itemize}
  \item $\J$ is the skew-symmetric structure matrix (symplectic part),
  \item $\mathbf{D}(\PhaseState) \succeq 0$ is the \emph{damping / dissipation} matrix
        (symmetric positive semi-definite),
  \item $\mathbf{B}(\PhaseState)$ is the input port matrix, and
  \item $u$ is an external control input or correction signal.
\end{itemize}
The output (observation) port is $y = \mathbf{B}^\top \nabla_\PhaseState H$.

\begin{definition}[Energy balance and passivity]
\label{def:passivity}
Along any trajectory of \eqref{eq:port_ham}, the stored energy $H$ satisfies
\begin{equation}
  \dot{H}
  = -(\nabla_\PhaseState H)^\top \mathbf{D}\,\nabla_\PhaseState H
    + y^\top u
  \;\leq\; y^\top u.
  \label{eq:energy_balance}
\end{equation}
The system is \textbf{passive} with storage function $H$ if
$\mathbf{D} \succeq 0$: energy stored is bounded above by energy
supplied through the port $y^\top u$.
Passivity is the central stability concept for our framework:
an adaptation step that preserves $\mathbf{D} \succeq 0$ is
guaranteed not to inject spurious energy into the dynamics.
\end{definition}

The GRL-SNAM Hamiltonian~\citep{ellendula2025grlsnam} (Definition~\ref{def:mech_hamiltonian}) is a
mechanical port-Hamiltonian where the dissipative port models motion
resistance and the input port receives online correction signals from
the adaptation loop.

Figure~\ref{fig:prelim:portham} illustrates the port-Hamiltonian structure.

\begin{figure}[t]
\centering
\begin{tikzpicture}[
  font=\small\sffamily,
  box/.style={rectangle, rounded corners=4pt,
              draw=#1!70!black, fill=#1!10,
              line width=1pt,
              minimum width=2.6cm, minimum height=1.1cm,
              align=center, inner sep=6pt},
  arr/.style={-{Stealth[length=6pt]}, line width=1.1pt, color=#1},
  lbl/.style={font=\scriptsize\sffamily\itshape, text=gray!60!black},
  eqlbl/.style={font=\scriptsize\sffamily, text=black},
]

% Central storage box
\node[box=teal] (H) at (0,0)
  {$H(\vec{q},p;\mathcal{E})$\\[2pt]\scriptsize stored energy};

% Symplectic flow
\node[box=blue!60!black] (flow) at (-4.2,0)
  {$\J\,\nabla_z H$\\[2pt]\scriptsize conservative flow\\(symplectic)};

% Dissipation
\node[box=red!60!black] (diss) at (0,-2.8)
  {$-\mathbf{D}\,\nabla_z H$\\[2pt]\scriptsize dissipation\\$\mathbf{D}\succeq 0$};

% External port
\node[box=orange] (port) at (4.2,0)
  {$\mathbf{B}u$\\[2pt]\scriptsize control /\\correction port};

% Output
\node[box=gray!60!black] (out) at (0,2.8)
  {$y = \mathbf{B}^\top\nabla_z H$\\[2pt]\scriptsize observation port};

% Summation
% \node[circle, draw=black!50, fill=white, minimum size=22pt,
%       font=\scriptsize\sffamily] (sum) at (0+0.1*0,0) {};
% \node at (sum) {$+$};

% Arrows
\draw[arr=blue!60!black] (flow.east) -- (H.west)
  node[midway, above, lbl] {$\dot{z}_{\mathrm{cons}}$};
\draw[arr=red!60!black]  (diss.north) -- (H.south)
  node[midway, right, lbl] {$\dot{z}_{\mathrm{diss}}$};
\draw[arr=orange]        (port.west) -- (H.east)
  node[midway, above, lbl] {$\dot{z}_{\mathrm{ctrl}}$};
\draw[arr=teal]          (H.west) -- (flow.east)
  node[midway, above=4pt, lbl] {$\nabla_z H$};
\draw[arr=teal]          (H.south) -- (diss.north);
\draw[arr=teal]          (H.east) -- (port.west);
\draw[arr=gray!60!black] (H.north) -- (out.south)
  node[midway, right, lbl] {$\dot{z}$};

% Energy balance annotation
\node[eqlbl, below=8pt of diss, text width=7cm, align=center]
  {$\dot{H} = -(\nabla_z H)^\top \mathbf{D}\,\nabla_z H + y^\top u
    \;\leq\; y^\top u$};

\end{tikzpicture}
\caption{%
  \textbf{Port-Hamiltonian structure.}
  The dynamics decompose into three additive contributions:
  a conservative symplectic flow (blue),
  a dissipative damping term (red, $\mathbf{D}\succeq 0$),
  and an external control/correction port (orange).
  The energy balance~\eqref{eq:energy_balance} ensures passivity:
  stored energy can only increase when energy is supplied through the port.
  Chapters~\ref{chap:grlsnam}--\ref{chap:material} use this structure to
  monitor stability of each adaptation step.
}
\label{fig:prelim:portham}
\end{figure}

\subsection{From Optimal Control to Reduced Hamiltonians}
\label{subsec:prelim:pmp}

At stage $\tau$ with context $\mathcal{E}_\tau$, navigation is posed as a
constrained optimal control problem (OCP):
\begin{equation}
\label{eq:ocp_prelim}
\begin{aligned}
\min_{\vec{u}(\cdot)\in\mathcal{U}}\;
&\int_0^T L(\vec{q}(t),\vec{u}(t);\mathcal{E}_\tau)\,dt \\
\text{s.t.}\quad
&\dot{\vec{q}} = f(\vec{q}) + A(\vec{q})\,\vec{u},\;
d_i(\vec{q};\mathcal{E}_\tau)\geq 0,\;
\vec{q}(0)=\vec{q}_0,\;\vec{q}(T)=\vec{q}_{\mathrm{target}}.
\end{aligned}
\end{equation}
The \emph{Pontryagin Maximum Principle} (PMP)~\citep{
FERREIRA1994whenpmp,hartl1995survey} provides necessary optimality conditions
by introducing a costate $p : [0,T]\to\mathbb{R}^n$ and defining the
\emph{control Hamiltonian}
$\mathcal{H}_c(\vec{q},p,\vec{u}) = \langle f(\vec{q})+A(\vec{q})\vec{u},\,p\rangle
- L(\vec{q},\vec{u};\mathcal{E})$.
The PMP motivates using a costate variable $p$ and a reduced phase-space
dynamics that jointly govern the optimal trajectory.

\textbf{Thesis surrogate.}
Rather than deriving the exact PMP Hamiltonian for an arbitrary navigation OCP
(which yields $H_{\mathrm{PMP}} = p^\top f(q) + \tfrac{1}{2}p^\top A R_u^{-1} A^\top p - \ell(q;\mathcal{E})$
under quadratic control cost, a sign convention unsuited to energy shaping),
we employ a \emph{mechanical surrogate Hamiltonian}:
\begin{equation}
H(\vec{q},p;\mathcal{E})
= \tfrac{1}{2}\,p^\top\Mass^{-1}p
  + V(\vec{q};\mathcal{E}),
\quad
V(\vec{q};\mathcal{E})
= \ell(\vec{q};\mathcal{E})
  + \textstyle\sum_{i=1}^{C(\mathcal{E})} \alpha_i(\mathcal{E})\,
    b_i\!\bigl(d_i(\vec{q};\mathcal{E})\bigr),
\label{eq:reduced_ham_prelim}
\end{equation}
where $b_i$ is a $C^2$ barrier potential, $\alpha_i(\mathcal{E})\geq 0$
are environment-conditioned weights, and the potential $V$ is chosen so
that its gradient acts as a repulsive/attractive force aligned with the
navigation objective.
This surrogate is \emph{inspired} by the PMP structure (a costate $p$ as
internal evaluator, Hamilton's equations as motion law) but is not claimed
to be the exact PMP Hamiltonian for arbitrary navigation OCPs.
It is learned end-to-end so that the induced flow \eqref{eq:policy_as_flow}
recovers navigation-optimal behavior.

% =======================================================================
\section{Stability, Barriers, and Safety}
\label{sec:prelim:stability}
% =======================================================================

\subsection{Lyapunov Stability and Energy-Based Analysis}
\label{subsec:prelim:lyapunov}

\begin{definition}[Lyapunov function]
\label{def:lyapunov}
A function $V : T^*\mathcal{Q} \to \mathbb{R}_{\geq 0}$ is a
\emph{Lyapunov function} for an equilibrium $\PhaseState^*$ if
$V(\PhaseState^*) = 0$, $V(\PhaseState) > 0$ for $\PhaseState \neq \PhaseState^*$,
and $\dot{V}(\PhaseState) \leq 0$ along trajectories.
Strict decrease ($\dot{V} < 0$) gives asymptotic stability.
\end{definition}

For a mechanical Hamiltonian \eqref{eq:mech_hamiltonian} with a single goal
at $\vec{q}^*$, the Hamiltonian itself is a Lyapunov candidate: dissipation
($\mathbf{D} \succ 0$) ensures $\dot{H} < 0$ when moving, driving the
system toward the energy minimum.
This is the stability guarantee used in Chapter~\ref{chap:grlsnam}.

\subsection{Control Barrier Functions}
\label{subsec:prelim:cbf}

\begin{definition}[Control Barrier Function (CBF)~\citep{ames2016control}]
\label{def:cbf}
Let $\mathcal{S} = \{\vec{q} : h(\vec{q}) \geq 0\}$ be a safe set.
A function $h \in C^1$ is a \emph{Control Barrier Function} if there exists
a class-$\mathcal{K}$ function $\kappa$ such that for all $\vec{q} \in \mathcal{S}$:
\begin{equation}
  \sup_{\vec{u}}\;\bigl[
    L_f h(\vec{q}) + L_g h(\vec{q})\,\vec{u}
  \bigr] \geq -\kappa(h(\vec{q})),
  \label{eq:cbf_condition}
\end{equation}
where $L_f h$ and $L_g h$ are Lie derivatives.
If \eqref{eq:cbf_condition} holds, then $\mathcal{S}$ is forward invariant.
\end{definition}

In our framework, the clearance functions $d_i(\vec{q};\mathcal{E})$
play the role of barrier certificates.
Rather than imposing them as external QP filters, they enter the energy as
barrier potentials $b_i(d_i)$, so that the gradient field naturally repels
trajectories from $\{d_i = 0\}$.
Chapter~\ref{chap:material} adds a steep hazard-barrier term
$H_{\mathrm{haz}}^{\mathrm{hard}}$ whose potential grows sharply near hazard
boundaries, creating a strong repulsive force.
Formal forward invariance under this barrier requires additional assumptions:
in continuous time, an infinite barrier with bounded energy suffices;
for finite-step rollouts, either a timestep bound or a discrete CBF
projection step is needed (see §\ref{subsec:ch4:dynamics}).

\subsection{Barrier Relaxation of State Constraints}
\label{subsec:prelim:barrier_relax}

For learning and differentiable rollout, the hard constraints
$d_i \geq 0$ in \eqref{eq:ocp_prelim} are replaced by
log-barrier or smooth-barrier potentials~\citep{NocedalWright2006NumericalOptimization}.

\begin{proposition}[Barrier relaxation, adapted from~\citep{malisani2023interior}]
\label{prop:barrier_convergence}
Under standard Slater regularity and second-order sufficiency conditions,
the solution of the barrier-relaxed OCP converges to the solution of the
state-constrained OCP as the barrier stiffness parameter $t \to \infty$.
The barrier-induced approximate multipliers
$\lambda_i^t(\tau) = \alpha_i/(t\,d_i(q^t(\tau)))$
converge weakly-$*$ to the true PMP constraint multipliers, and
complementarity holds in the limit.
\end{proposition}

This proposition licenses the use of differentiable barrier potentials in
the learned Hamiltonian: the training loss can be backpropagated through
the rollout, and as the learned $\alpha_i$ values grow, the relaxed solution
converges to the constrained one.

% =======================================================================
\section{Risk-Sensitive Objectives}
\label{sec:prelim:risk}
% =======================================================================

Average-cost optimization is insufficient when the environment contains
rare but catastrophic material hazards: a policy that maximizes expected
return may still exhibit high tail-risk.

\begin{definition}[Value-at-Risk and CVaR]
\label{def:cvar}
For a random cost $X$ and confidence level $\beta \in (0,1)$:
\begin{align}
  \mathrm{VaR}_\beta(X) &:= \inf\{x : P(X \leq x) \geq \beta\},
  \label{eq:var} \\
  \mathrm{CVaR}_\beta(X) &:= \min_{\nu \in \mathbb{R}}
    \left[\nu + \frac{1}{1-\beta}\,\mathbb{E}(X - \nu)_+\right].
  \label{eq:cvar}
\end{align}
For continuous distributions this simplifies to
$\mathrm{CVaR}_\beta(X) = \mathbb{E}[X \mid X \geq \mathrm{VaR}_\beta(X)]$.
$\mathrm{CVaR}_\beta$ is the expected cost in the upper $(1-\beta)$ tail of
outcomes: for $\beta = 0.95$ it focuses on the worst $5\%$ of realizations,
not the worst $95\%$.
It is a \emph{coherent risk measure}~\citep{rockafellar2002cvar}:
convex, monotone, and sub-additive, making it tractable as an optimization
objective.
\end{definition}

\begin{definition}[Constrained MDP (CMDP)~\citep{achiam2017constrained}]
\label{def:cmdp}
A CMDP augments a standard MDP with a cost function $c : \mathcal{S}\times\mathcal{A}\to\mathbb{R}$
and a constraint $\mathbb{E}\bigl[\sum_t \gamma^t c(s_t,a_t)\bigr] \leq \kappa$.
Constrained policy optimization (CPO) solves CMDPs with near-constraint
satisfaction per update step.
\end{definition}

In Chapter~\ref{chap:material}, the tail-risk objective enters as a
CVaR term over material cost along the trajectory:
\begin{equation}
  J_{\mathrm{risk}}(\pi)
  = \mathrm{CVaR}_\beta\!\Bigl[\int_0^T \rho(\vec{c}(t))\,dt\Bigr],
  \label{eq:cvar_objective}
\end{equation}
where $\rho : \mathcal{Q}_f \to [0,1]$ is the material risk field.
This is not a constraint added externally: it induces a new energy factor
$H_{\mathrm{mat}}^{\mathrm{soft}}$ whose gradient directly reshapes the
force field.

% =======================================================================
\section{Multi-Timescale Adaptation Hierarchy}
\label{sec:prelim:continual}
% =======================================================================

A defining feature of our approach is that adaptation operates at
\emph{three nested timescales}, each correcting a different aspect of the
agent's behavior.
Figure~\ref{fig:prelim:timescales} illustrates this hierarchy.

\begin{figure}[htbp]
\centering
\begin{tikzpicture}[
  font=\small\sffamily,
  ring/.style args={#1,#2,#3}{
    ellipse, draw=#1!70!black, fill=#1!8,
    line width=1.2pt,
    minimum width=#2, minimum height=#3,
    align=center
  },
  lbl/.style args={#1}{
    font=\scriptsize\sffamily, text=#1
  },
  arr/.style={-{Stealth[length=5pt]}, line width=1pt, dashed, color=gray!60},
]

% Three nested ellipses
\node[ring={violet,6.4cm,3.4cm}] (outer) at (0,0) {};
\node[ring={orange,4.4cm,2.3cm}] (mid)   at (0,0) {};
\node[ring={teal,2.4cm,1.2cm}] (inner)   at (0,0) {};

% Labels inside
\node[lbl={teal!80!black}, font=\small\sffamily\bfseries] at (0,0)
  {Step loop};
\node[lbl={teal!60!black}] at (0,-0.38)
  {$\Phi^{\Delta t}_H$ integration};

\node[lbl={orange!80!black}, font=\small\sffamily\bfseries]
  at (0, 1.32) {Segment / episode loop};
\node[lbl={orange!60!black}]
  at (0, 0.82) {online $\alpha_i, \Mass$ updates};

\node[lbl={violet!80!black}, font=\small\sffamily\bfseries]
  at (0, 2.18) {Curriculum loop};
\node[lbl={violet!60!black}]
  at (0, 1.86) {meta-parameter $\xi$, coverage, risk};

% Side annotations
\node[lbl={teal!70!black}, right=2.5cm of inner,
      text width=3.2cm, align=left]
  (ann1) {$\Delta t \sim 10$ms\\Symplectic integration\\Fixed $H$ per stage};

\node[lbl={orange!70!black}, above right=-1.2cm and -9.0cm of mid,
      text width=3.4cm, align=left]
  (ann2) {$\sim$episode / corridor\\Update barrier weights $\alpha_i$\\Online correction signals};

\node[lbl={violet!70!black}, above right=0.8cm and 0.4cm of outer,
      text width=3.8cm, align=left]
  (ann3) {$\sim$curriculum / deployment\\Update $\xi$: energy parameterization\\Risk model, coverage objectives};

% Arrows
\draw[arr] (ann1.west) -- +(-2.6,0);
\draw[arr] (ann2.east) -- +(1.4,0);
\draw[arr] (ann3.south west) -- +(-0.3,-0.6);

\end{tikzpicture}
\caption{%
  \textbf{Multi-timescale adaptation hierarchy.}
  Three nested loops operate at increasing timescales.
  The \emph{step loop} (innermost, teal) integrates the fixed Hamiltonian
  $H$ via a symplectic solver at each $\Delta t$.
  The \emph{segment/episode loop} (middle, orange) updates barrier weights
  $\alpha_i(\mathcal{E})$ and correction signals online, reshaping the
  energy landscape between corridor segments.
  The \emph{curriculum loop} (outer, violet) updates the meta-parameters
  $\xi$ of the scenario-to-energy network, incorporating risk objectives
  and coverage criteria over full training episodes.
  Each loop corrects a different aspect of the decision field
  without disrupting the inner loops' structure.
}
\label{fig:prelim:timescales}
\end{figure}

\paragraph{Step loop.}
At every integration step $t \to t + \Delta t$, the symplectic integrator
advances the phase-space state using the current Hamiltonian $H(\cdot;\mathcal{E}_\tau)$.
No parameters are updated here; the field topology is fixed for the stage.

\paragraph{Segment / episode loop.}
After each navigation segment or episode, the barrier weights
$\{\alpha_i(\mathcal{E}_\tau)\}$ and online correction signals $u$ are updated
based on clearance violations, goal proximity, and (in Stage 2) material-risk
observations.
This is the \emph{geometric alignment} step that reshapes the energy landscape
between stages without catastrophic relearning.

\paragraph{Curriculum loop.}
Over longer deployment, the meta-parameters $\xi$ of the
scenario-to-energy mapping $\eta_\xi : \mathcal{E}_\tau \mapsto H(\cdot;\mathcal{E}_\tau)$
are updated using rollout gradients, risk objectives, and coverage criteria.
This loop corresponds to the continual learning component of the framework,
and is the focus of Chapter~\ref{chap:material}.

% =======================================================================
\section{Learning from Demonstrations and Rewards}
\label{sec:rw:demonstrations}
% =======================================================================

\paragraph{Imitation learning.}
Behavioral cloning~\citep{torabi2018behavioralcloningobservation} treats policy learning as
supervised regression onto expert actions, scaling well with data but
degrading under distribution shift.
DAgger~\citep{shi2025dagger} corrects this by querying the expert
on learner-visited states, and GAIL~\citep{ho2016generativeadversarialimitationlearning} reframes
imitation as occupancy-measure matching.
For deformable settings, IL has been applied to manipulation
trajectories~\citep{zhang2024adaptigraph} and coordinated
motion~\citep{Lai_2025}.
\emph{Gap:} all IL variants fix the demonstration target and learn a
policy mapping; the resulting controller exposes no force-field
topology and has no gradient pathway for reshaping dynamics when
the environment changes.

\paragraph{Reinforcement learning for navigation.}
PPO~\citep{schulman2017proximalpolicyoptimizationalgorithms}, SAC~\citep{haarnoja2018softactorcriticoffpolicymaximum,mohammad2025sac},
and TD3~\citep{fujimoto2018addressingfunctionapproximationerror} dominate continuous-control navigation.
Large-scale platforms (Habitat, Gibson, HM3D) and distributed
training~\citep{wijmans2019dd} have pushed PointNav near-perfect in
simulation, but performance degrades sharply under distribution shift.
Hierarchical RL~\citep{vezhnevets2017feudalnetworkshierarchicalreinforcement} adds temporal
structure without exposing the topology of the closed-loop field.
Most methods remain within the dynamic programming principle; a
growing line of work departing from it like HJB-based
learning~\citep{settai2025temporal}, martingale
formulations~\citep{jia2023qlearning}, and duality-driven Hamiltonian
RL~\citep{nguyen2024dpo} is more directly related to this thesis.
\emph{Gap:} the reward is fixed; adapting to new material or risk
requires reward redesign, not field reshaping.

\paragraph{Inverse RL and reward learning.}
MaxEnt IRL~\citep{ziebart2008maximum} and adversarial
IRL~\citep{luo2024c} partially relax the fixed-reward assumption
by recovering objectives from demonstrations.
Terrain-cost learning~\citep{sathyamoorthy2022terrapn} recovers
traversability fields from demonstrations for off-road navigation.
\emph{Gap:} IRL learns the reward but fixes the downstream optimizer;
updates to the objective do not automatically propagate as stable
updates to the dynamics, and the closed-loop topology remains
unconstrained.

% =======================================================================
\section{Structured Control and Planning under Constraints}
\label{sec:rw:structured}
% =======================================================================

\paragraph{Classical planning and model-based navigation.}
A*~\citep{hart1968formal} and sampling-based planners
(RRT*, PRM*)~\citep{karaman2011sampling} provide completeness and
optimality guarantees over explicit cost fields.
MPPI~\citep{williams2017mppi} enables stochastic optimal control
with GPU-parallel rollouts.
Neural scene representations like NeRF-SLAM~\citep{Zhu2022CVPR,sucar2021imap},
GS-SLAM~\citep{keetha2024splatamsplattrack} and differentiable
planners~\citep{gupta2019cognitivemappingplanningvisual} further enrich the map for
downstream use.
\emph{Gap:} planners fix the cost/objective representation and optimize
paths relative to it; the decision field structure is not itself learned
or continually updated as latent semantic properties change.

\paragraph{Safety-critical control.}
CBF-based QP filters~\citep{ames2016control} enforce forward
invariance of a safe set at runtime; neural zeroing
barriers~\citep{feng2023safergapgapbasedlocal} and adaptive
mechanisms~\citep{mohammad2025sac} extend these to learned settings.
CPO~\citep{achiam2017constrained} provides near-constraint satisfaction
per gradient step under CMDPs.
CVaR MDPs~\citep{chow2015risk} optimize tail-risk rather than expected
cost, focusing on rare catastrophic outcomes.
\emph{Gap:} safety and risk typically appear as external projection
layers or constraint penalties on top of unstructured controllers.
We integrate them directly as energy terms whose gradients
shape the force field, co-optimizing safety and task performance under
a single structure-preserving flow.

\paragraph{Geometric, Hamiltonian, and port-Hamiltonian learning.}
HNNs~\citep{greydanus2019hamiltonian} learn a scalar $H(q,p)$ whose
gradients define conservative dynamics, improving long-horizon
stability over unconstrained neural ODEs~\citep{chen2019neuralordinarydifferentialequations}.
SymODEN~\citep{zhong2020symplectic} adds control inputs; Lagrangian
networks~\citep{cranmer2020lagrangian} remove the need for canonical
coordinates.
Port-Hamiltonian networks~\citep{dipersio2025porthamiltonianneuralnetworkstheory,desai2021porthamiltonian}
explicitly separate stored energy, dissipation, and forcing ports;
stable port-HNNs~\citep{roth2025stable} enforce global Lyapunov
stability by construction.
Geometric RL with equivariant
policies~\citep{hoang2025geometry,martínezrubio2023acceleratedriemannianoptimizationhandling} and
duality-driven Hamiltonian RL~\citep{nguyen2024dpo} exploit
configuration-space structure for manipulation and optimal control
without Bellman backups.
Symplectic integrators~\citep{reich1996symplectic,hairer2006geometric}
preserve the symplectic form over long rollouts, controlling the
energy drift that plagues generic integrators.
Table~\ref{tab:rw:hamiltonian} summarizes this literature.
\emph{Gap:} existing Hamiltonian and port-Hamiltonian learning is
developed for \emph{system identification} for recovering the energy of a
known physical system from trajectory data and not for continual
risk-sensitive navigation in unknown environments.
We transfer the same structural principles to decision-making:
energy factors encode navigation objectives, and adaptation loops
update coefficients online as new structure is revealed.

\begin{table}[t]
\centering
\caption{%
  \textbf{Geometric and Hamiltonian learning literature.}
  No existing method provides navigation decision-making with online
  adaptation \emph{and} risk-semantic terms simultaneously.
}
\label{tab:rw:hamiltonian}
\small
\setlength{\tabcolsep}{5pt}
\renewcommand{\arraystretch}{1.3}
\resizebox{\textwidth}{!}{
\begin{tabular}{lcccc}
\toprule
\textbf{Method} &
\makecell{\textbf{Non-conserv.}\\\textbf{/ control}} &
\makecell{\textbf{Navigation}\\\textbf{decision}} &
\makecell{\textbf{Online}\\\textbf{adapt.}} &
\makecell{\textbf{Risk /}\\\textbf{semantic}} \\
\midrule
HNN~\citep{greydanus2019hamiltonian}            & \none & \none & \none & \none \\
SymODEN~\citep{zhong2020symplectic}             & \partially & \none & \none & \none \\
LNN~\citep{cranmer2020lagrangian}               & \none & \none & \none & \none \\
Port-HNN~\citep{desai2021porthamiltonian} & \full & \none & \none & \none \\
Stable port-HNN~\citep{roth2025stable}          & \full & \none & \none & \none \\
Geometric RL~\citep{hoang2025geometry}          & \none & \partially & \none & \none \\
Hamiltonian RL~\citep{nguyen2024dpo}            & \partially & \full & \none & \none \\
HJB-based RL~\citep{settai2025temporal}         & \partially & \full & \none & \none \\
\midrule
\textbf{Stage~1 (Ch.~\ref{chap:grlsnam})}  & \full & \full & \full & \none \\
\textbf{Stage~2 (Ch.~\ref{chap:material})} & \full & \full & \full & \full \\
\bottomrule
\end{tabular}}
\end{table}

% =======================================================================
\section{Learning and Adapting in Unknown Environments}
\label{sec:rw:unknown}
% =======================================================================

\paragraph{Semantic and material-aware navigation.}
TerraPN~\citep{sathyamoorthy2022terrapn} learns terrain cost maps from
demonstrations for a DWA~\citep{fox1997dynamic} downstream planner.
Context-aware reward learning~\citep{wang2024crl} adapts objectives
to semantic context.
Neural Topological SLAM~\citep{chaplot2020neuraltopologicalslamvisual} and semantic vision
extensions~\citep{zhang2025semantic} provide richer map representations
for goal-directed navigation.
RUGD~\citep{rugd2019} and RELLIS-3D~\citep{rellis2020} supply off-road
semantic datasets; self-supervised traversability
methods~\citep{mclennan2025learning} recover material costs without
annotation.
\emph{Gap:} these methods treat perception as upstream and planning as
downstream; semantic signals update the map but not the structure of the
decision dynamics.
We integrate material signals as explicit energy terms whose
gradients directly reshape the phase-space vector field.

\paragraph{Simultaneous navigation and mapping (SNAM).}
SGoLAM~\citep{kim2021sgolamsimultaneousgoallocalization}, CMP~\citep{gupta2019cognitivemappingplanningvisual}, and
Active Neural SLAM~\citep{chaplot2020learning} couple map construction
with policy execution, prioritizing rich global maps.
They do not reason about mapping cost and are not formulated as optimal
control problems.
GRL-SNAM instead targets joint exploration with minimal sensing effort:
the energy landscape is refined only in the local neighborhood the agent
visits, not rebuilt globally.

\paragraph{Continual learning, meta-learning, and online adaptation.}
EWC~\citep{kirkpatrick2017ewc} protects parameters from forgetting via
Fisher-information regularization; CLEAR~\citep{rolnick2019experience}
adds experience replay; progressive networks~\citep{rusu2016progressive}
expand architecture per task.
MAML~\citep{finn2017maml} learns initializations for rapid fine-tuning;
RMA~\citep{kumar2021rma} adds an online adaptation module for terrain
and payload changes in legged locomotion.
\emph{Gap:} CL and meta-learning adapt parameters and latent context but
do not regulate the topology or geometry of the closed-loop vector
field.
We contribute a complementary axis: continual reshaping of
the energy landscape itself, monitored by passivity and coverage
criteria across three nested timescales.

\paragraph{Foundation models and LLM-based planning.}
SayCan~\citep{ahn2022saycan}, RT-1/RT-2~\citep{brohan2023rt1roboticstransformerrealworld,rt22023},
Code as Policies~\citep{liang2022code}, and LLM-based
navigation~\citep{zhu2024navi2gazeleveragingfoundationmodels,xiong2026sfconavefficientzeroshotvisual}
provide strong semantic goal specification and task decomposition.
They govern the high-level planning layer and rely on separate
low-level controllers for stable continuous dynamics.
Our approach is complementary: LLMs can supply semantic goals encoded
as energy-term coefficients; the Hamiltonian dynamics then ensures
physically stable, risk-aware execution.

Foundation models and the Hamiltonian navigation framework are
complementary, not competing.
LLMs can supply semantic goals, task decompositions, and material-context
priors that enter our framework as potential energy terms; the
Hamiltonian dynamics then ensures that these high-level intentions translate
into physically stable, risk-aware trajectories.

% -----------------------------------------------------------------------
\section{Summary and Positioning of This Thesis}
\label{sec:rw:summary}
% -----------------------------------------------------------------------

Figure~\ref{fig:rw:gapmap} synthesizes the paradigm landscape along two axes:
the degree of \emph{structural dynamics} encoded in the policy, and the
\emph{timescale of adaptation}.
Three observations emerge.

First, methods with the strongest structural guarantees (port-Hamiltonian
networks, CBF-based controllers, classical planners) tend to fix either the
dynamics class or the objective, and provide little or no continual
adaptation.

Second, methods with the richest continual adaptation (meta-RL, CL,
LLM-planners) typically operate over flat policy or parameter spaces, with no
explicit topology guarantees on the resulting dynamics.

Third, the combination of structured dynamics and continual adaptation, the
combination needed for safe, correctable, risk-sensitive navigation in unknown
environments, is the gap that our work directly addresses.

Figure~\ref{fig:rw:gapmap} motivates
our central contribution: a Hamiltonian learning framework in which
the policy is the flow map of an energy landscape that is structured by
design and continually enriched by experience, so that geometry, material
risk, and adaptation timescale are not separate concerns but unified aspects
of one framework.


\begin{figure}[t]
\centering
\begin{tikzpicture}[
  font=\small\sffamily,
  % styles
  blk/.style={
    rectangle, rounded corners=4pt,
    draw=#1!70!black, fill=#1!9,
    line width=1pt,
    minimum width=2.8cm, minimum height=0.9cm,
    align=center, inner sep=4pt
  },
  thesisblk/.style={
    rectangle, rounded corners=4pt,
    draw=violet!70!black, fill=violet!15,
    line width=1.5pt,
    minimum width=3.2cm, minimum height=0.9cm,
    align=center, inner sep=4pt,
    font=\small\sffamily\bfseries
  },
  gap/.style={
    draw=red!50!black, fill=red!6,
    dashed, line width=0.8pt
  },
  arr/.style={-{Stealth[length=5pt]}, line width=1pt, color=#1},
  lbl/.style={font=\scriptsize\sffamily\itshape, text=gray!70!black},
]

% ---- x-axis: "what is learned" (structure → flat) ----
% ---- y-axis: "adaptation timescale" (none → continual) ----

% Column headers
\node[font=\scriptsize\sffamily\bfseries, text=gray!60] at (-3.8, 3.6)
  {Structured dynamics};
\node[font=\scriptsize\sffamily\bfseries, text=gray!60] at (3.8, 3.6)
  {Flat policy / value fn.};
\draw[arr=gray!40, <->] (-2.2,3.6) -- (2.2,3.6);

% Row labels
\node[font=\scriptsize\sffamily\bfseries, text=gray!60, rotate=90] at (-7.2, 0)
  {No adaptation $\longrightarrow$ Continual adaptation};
\draw[arr=gray!40, <->] (-6.8,-2.6) -- (-6.8, 2.6);

% ---- Paradigm boxes ----
% top-left: structured + no/partial adaptation
\node[blk=teal]   (hnn)   at (-3.6, 2.0) {Port-HNN\\(sysid only)};
\node[blk=teal]   (cbf)   at (-3.6, 0.6) {CBF / Safe RL\\(fixed safety set)};
\node[blk=teal]   (plan)  at (-3.6,-0.8) {A* / RRT* / MPC\\(fixed cost field)};

% top-right: flat + no/partial adaptation
\node[blk=orange] (bc)    at ( 3.6, 2.0) {BC / DAgger\\(fixed demos)};
\node[blk=orange] (rl)    at ( 3.6, 0.4) {PPO / SAC / TD3\\(fixed reward)};
\node[blk=orange] (irl)   at ( 3.6,-1.2) {IRL / GAIL\\(two-stage)};

% bottom-right: flat + continual adaptation
\node[blk=blue!60!black] (meta)  at ( 3.6,-2.9) {MAML / CLEAR\\(param.\ adaptation)};
\node[blk=blue!60!black] (llm)   at ( 0.0,-2.2) {LLM planners\\(semantic only)};

% bottom-left: structured + continual — THE GAP, then filled by thesis
\node[blk=red!60!black, gap] (gapbox) at (-3.6,-2.9)
  {\textcolor{red!70!black}{\textbf{GAP:}}\\ structured \&\\ continual};

% thesis arrow
\node[thesisblk] (thesis) at (-3.6, -5.0)
  {\textbf{This thesis}\\\textbf{(Continual Ham.\ Learning)}};
\draw[arr=violet!70!black, line width=1.5pt]
  (thesis.north) -- (gapbox.south)
  node[midway, right=2pt, font=\scriptsize\sffamily\itshape, text=violet!80!black]
    {fills this gap};

% arrows showing "each paradigm leaves something fixed"
\node[lbl, above=2pt of hnn]  {learns $H$ but not for navigation};
\node[lbl, above=2pt of bc]   {demo target is fixed};
\node[lbl, above=-2pt of rl]   {reward is fixed};
\node[lbl, below=2pt of plan] {cost field is fixed};
\node[lbl, below=2pt of irl]  {optimizer is fixed};
\node[lbl, below=2pt of meta] {controller form is fixed};

\end{tikzpicture}
\caption{%
  \textbf{Paradigm gap map.}
  Existing methods cluster into two axes: those providing structured dynamics
  (left column) and those with richer adaptation (right/bottom).
  No existing paradigm simultaneously provides structured dynamics \emph{and}
  continual adaptation with risk sensitivity.
  This thesis fills that gap through continual Hamiltonian learning.
}
\label{fig:rw:gapmap}
\end{figure}

% =======================================================================
\section{Supporting Theoretical Results}
\label{sec:prelim:theory}
% =======================================================================

The following lemmas and theorems underpin the claims made in
Chapters~\ref{chap:grlsnam} and~\ref{chap:material}.
They are stated precisely here and referenced throughout the thesis.

% -----------------------------------------------------------------------
\subsection*{A. Well-Posedness of the Learned Hamiltonian Flow}
% -----------------------------------------------------------------------

\begin{lemma}[Local existence and uniqueness]
\label{lem:wellposed}
Fix an active obstacle set $\mathcal{C}$, a positive-definite mass matrix
$\Mass \succ 0$, and bounded energy coefficients.
Suppose $H_\theta \in C^2$ on the safe interior
$\{\vec{q} : d_i(\vec{q}) > 0,\, \forall i \in \mathcal{C}\}$.
Then the vector field
\[
  \dot{z} = (\J - \mathbf{D}_\theta(z))\,\nabla_z H_\theta(z) + \mathbf{B}\,u
\]
is locally Lipschitz on the safe interior, so the rollout has a unique
solution on $[0,T]$ up to the first active-set change or boundary contact.
\end{lemma}

\begin{proof}[Proof sketch]
$\nabla_z H_\theta$ is $C^1$ (since $H_\theta \in C^2$) on the open safe
interior, and $\J$, $\mathbf{D}_\theta$, $\mathbf{B}$ are bounded there.
The Picard--Lindel\"of theorem then gives local existence and uniqueness.
\end{proof}

% -----------------------------------------------------------------------
\subsection*{B. Continuous-Time Passivity}
% -----------------------------------------------------------------------

\begin{theorem}[Port-Hamiltonian energy balance]
\label{thm:passivity}
If $\mathbf{D}_\theta(z) \succeq 0$ for all $z$, then along solutions of
\eqref{eq:port_ham}:
\[
  \dot{H}_\theta = -(\nabla_z H_\theta)^\top \mathbf{D}_\theta\,\nabla_z H_\theta
                  + y^\top u,
  \qquad y = \mathbf{B}^\top \nabla_z H_\theta.
\]
In particular, if $u = 0$ the stored energy is non-increasing:
$\dot{H}_\theta \leq 0$.
\end{theorem}

\begin{proof}
Differentiate $H_\theta(z(t))$ along trajectories and substitute
$\dot{z} = (\J - \mathbf{D}_\theta)\nabla_z H_\theta + \mathbf{B}u$.
Skew-symmetry of $\J$ gives $(\nabla_z H_\theta)^\top \J\,\nabla_z H_\theta = 0$.
\end{proof}

% -----------------------------------------------------------------------
\subsection*{C. Discrete Passivity under the Split Integrator}
% -----------------------------------------------------------------------

\begin{theorem}[Discrete energy inequality]
\label{thm:discrete_passivity}
Let the integration step split into a symplectic conservative half-step
(leapfrog) followed by a dissipative correction and port input.
For step size $\Delta t \leq h_0$ (a problem-dependent constant),
\[
  H_\theta(z_{n+1}) - H_\theta(z_n)
  \;\leq\;
  \Delta t\, y_n^\top u_n + C\,\Delta t^2,
\]
where $C$ depends on the Lipschitz constant of $\nabla_z H_\theta$ and the
magnitude of $\mathbf{D}_\theta$.
\end{theorem}

\begin{proof}[Proof sketch]
The conservative substep preserves $H_\theta$ up to $O(\Delta t^2)$
(symplectic integrator shadow Hamiltonian bound).
The dissipative substep decreases $H_\theta$ by at least
$\Delta t \,(\nabla_z H_\theta)^\top \mathbf{D}_\theta\,\nabla_z H_\theta \geq 0$.
The port-input substep adds $\Delta t\, y_n^\top u_n + O(\Delta t^2)$.
Combining gives the stated inequality.
\end{proof}

% -----------------------------------------------------------------------
\subsection*{D. Energy-Barrier Forward Invariance (Continuous Time)}
% -----------------------------------------------------------------------

\begin{theorem}[Barrier safety]
\label{thm:barrier_safety}
Let $\mathcal{S} = \{\vec{q} : d_i(\vec{q}) > 0,\, \forall i\}$ and
suppose $b_i(s) \to +\infty$ as $s \downarrow 0$.
If $H_\theta(z(0)) < B < \infty$ and the port input is energy-bounded so
that $H_\theta(z(t)) \leq B$ for all $t$ (e.g., $u = 0$ or passivity
ensures energy is non-increasing), then $\vec{q}(t) \in \mathcal{S}$ for
all $t \geq 0$ in the continuous-time system.

\emph{Discrete caveat:} For finite-step integration, this guarantee requires
either (i) a sufficiently small timestep $\Delta t \leq h_0$ such that
barrier overflow cannot occur in one step, or (ii) a CBF-based projection
\[
  h(\vec{q}_{n+1}) - h(\vec{q}_n) \geq -\alpha\, h(\vec{q}_n)
\]
enforced as a safety filter.
\end{theorem}

% -----------------------------------------------------------------------
\subsection*{E. Local Residual Decrease of Online Coefficient Update}
% -----------------------------------------------------------------------

\begin{lemma}[Tikhonov--Gauss--Newton step]
\label{lem:tikhonov}
The online coefficient update
$\Delta\zeta = (J_t^\top J_t + \lambda I)^{-1} J_t^\top r_t$
minimizes $\min_\Delta \|J_t \Delta - r_t\|^2 + \lambda\|\Delta\|^2$.
\end{lemma}

\begin{theorem}[Local observable residual decrease]
\label{thm:residual_decrease}
Suppose the observable map $y(\zeta)$ is $C^1$, and $J_t$ approximates
$\partial y/\partial\zeta$ with error $\epsilon_J$.
For step size $\kappa > 0$ sufficiently small,
\[
  \|y(\zeta + \kappa\Delta\zeta) - y^\star\|
  \;\leq\;
  \|y(\zeta) - y^\star\| - c\kappa\|J_t^\top r_t\|^2
  + O(\kappa^2 + \epsilon_J),
\]
where $c > 0$ depends on the regularization $\lambda$ and the local
Lipschitz constant of $y$.
\end{theorem}

% -----------------------------------------------------------------------
\subsection*{F. CVaR Gradient and Sample Complexity}
% -----------------------------------------------------------------------

\begin{theorem}[Empirical CVaR subgradient]
\label{thm:cvar_gradient}
For bounded episode cost $J_\theta \in [0, J_{\max}]$:
\[
  \operatorname{CVaR}_\beta(J_\theta)
  = \min_{\nu \in \mathbb{R}}
    \Bigl[\nu + \tfrac{1}{1-\beta}\,\mathbb{E}(J_\theta - \nu)_+\Bigr].
\]
A subgradient with respect to $\theta$ at the optimal $\nu^\star$ is:
\[
  \nabla_\theta \operatorname{CVaR}_\beta(J_\theta)
  = \frac{1}{1-\beta}\,
    \mathbb{E}\!\Bigl[\mathbf{1}\{J_\theta \geq \nu^\star\}\,
    \nabla_\theta J_\theta \Bigr].
\]
\end{theorem}

\begin{lemma}[CVaR concentration]
\label{lem:cvar_concentration}
If $J \in [0, J_{\max}]$, then with $B$ i.i.d.\ episode samples,
\[
  \bigl|\widehat{\operatorname{CVaR}}_\beta - \operatorname{CVaR}_\beta\bigr|
  = O\!\left(\frac{J_{\max}}{1-\beta}\sqrt{\frac{\log(1/\delta)}{B}}\right)
\]
with probability $1 - \delta$.
This bounds the number of episodes required for reliable CVaR updates.
\end{lemma}

% -----------------------------------------------------------------------
\subsection*{G. Coefficient Identifiability}
% -----------------------------------------------------------------------

\begin{theorem}[Local coefficient identifiability]
\label{thm:identifiability}
Let the force channels be $\phi_i(\vec{q}) = -\nabla_{\vec{q}} H_i(\vec{q})$.
If the Gram matrix
\[
  G_{ij} = \sum_t \phi_i(\vec{q}_t)^\top \phi_j(\vec{q}_t)
\]
is full rank over the training trajectories, then the coefficient vector
$(\beta, \lambda_s, \lambda_h, \{\alpha_i\}, \gamma)$ is locally
identifiable from trajectory data up to measurement noise.
\end{theorem}

% -----------------------------------------------------------------------
\subsection*{H. Per-Step Computational Complexity}
% -----------------------------------------------------------------------

\begin{proposition}[Per-step complexity]
\label{prop:complexity}
For active set size $m = |\mathcal{C}_t|$, state dimension $n$, and
coefficient dimension $k = m + O(1)$, one integration/update step costs:
\[
  O(mn) \quad \text{(force evaluation)},
  \qquad
  O(k^3) \quad \text{(Tikhonov solve, or } O(k^2) \text{ with rank-1 update)}.
\]
This is local in the active set $m \ll |\mathcal{W}|$ and does not scale
with the global workspace size, supporting efficiency claims in §\ref{sec:ch3:efficiency}.
\end{proposition}
